% GENERATED FILE. Edit canonical lessons, then run npm run generate:books.
% canonical-source-hash: fnv1a64:ee4d4970d660953b
% canonical-lessons: TA-C181-pitittamana, TA-C181-mocamana, TA-C181-pala, TA-C181-ellam, TA-C181-tatimanana

\chapter{Favourite, Terrible, Many, Everything, Thick}
\label{ch:ta-favourite-terrible-many-everything-thick}
\hlchaptermodality{181}

\begin{chapteropening}
\textbf{By the end of this chapter:} \emph{I can say favourite, terrible, many, everything and thick.}

Complete the last lesson of chapter 181: I can say favourite, terrible, many, everything and thick.
\end{chapteropening}

\section[\texorpdfstring{\ta{பிடித்தமான} (piṭittamāṉa)}{பிடித்தமான (piṭittamāṉa)}]{\ta{பிடித்தமான} (piṭittamāṉa) — favourite}

\label{lesson:TA-C181-pitittamana}



\begin{quote}
Before the new one: say the Tamil for funny, then the Tamil for interesting.
\end{quote}

\subsection*{You'll want to know: \ta{பிடித்தமான}}
\textbf{\ta{பிடித்தமான}} — \emph{piṭittamāṉa} — \textquotedblleft{}favourite\textquotedblright{}.

\begin{cousinweb}
\textbf{\ta{பிடித்தமான}} (\emph{piṭittamāṉa}) is native Dravidian (DEDR 4148). The entry glosses the verb \emph{piṭi} \textquotedblleft{}to grasp\textquotedblright{} and also \textquotedblleft{}to be pleasing\textquotedblright{}; what you are fond of is what \textquotedblleft{}takes hold\textquotedblright{} of you, as in \emph{enakku piṭikkum}, \textquotedblleft{}I like it\textquotedblright{}.
\end{cousinweb}

\begin{grammarlens}[title={What you've built}]
One more word for this chapter.
\end{grammarlens}

\subsection*{Your turn}
\begin{itemize}
\raggedright
  \item \emph{Say it:} \emph{piṭittamāṉa}
  \item \emph{Say it:} \emph{piṭittamāṉa}, once more
  \item \emph{Say it:} say \emph{piṭittamāṉa}
  \item \emph{Recall:} say the Tamil for funny, then the Tamil for interesting, then say \emph{piṭittamāṉa} again
\end{itemize}

\begin{tcolorbox}[breakable,colback=teal!4,colframe=teal!35!black,title={Before you move on}]
What does \ta{பிடித்தமான} mean? (\textquotedblleft{}Favourite\textquotedblright{}.) Say it once more.
\end{tcolorbox}
\section[\texorpdfstring{\ta{மோசமான} (mōcamāṉa)}{மோசமான (mōcamāṉa)}]{\ta{மோசமான} (mōcamāṉa) — terrible, awful}

\label{lesson:TA-C181-mocamana}



\begin{quote}
Before the new one: say the Tamil for interesting, then the Tamil for favourite.
\end{quote}

\subsection*{You'll want to know: \ta{மோசமான}}
\textbf{\ta{மோசமான}} — \emph{mōcamāṉa} — \textquotedblleft{}terrible, awful\textquotedblright{}.

\begin{cousinweb}
The origin of \textbf{\ta{மோசமான}} (\emph{mōcamāṉa}) is uncertain. \emph{Mōcam}, \textquotedblleft{}deceit, a bad turn\textquotedblright{}, has no entry in DEDR; it is probably Sanskrit \emph{mōṣa}, \textquotedblleft{}theft\textquotedblright{} (CDIAL 10358).
\end{cousinweb}

\begin{grammarlens}[title={What you've built}]
One more word for this chapter.
\end{grammarlens}

\subsection*{Your turn}
\begin{itemize}
\raggedright
  \item \emph{Say it:} \emph{mōcamāṉa}
  \item \emph{Say it:} \emph{mōcamāṉa}, once more
  \item \emph{Say it:} say \emph{mōcamāṉa}
  \item \emph{Recall:} say the Tamil for interesting, then the Tamil for favourite, then say \emph{mōcamāṉa} again
\end{itemize}

\begin{tcolorbox}[breakable,colback=teal!4,colframe=teal!35!black,title={Before you move on}]
What does \ta{மோசமான} mean? (\textquotedblleft{}Terrible, awful\textquotedblright{}.) Say it once more.
\end{tcolorbox}
\section[\texorpdfstring{\ta{பல} (pala)}{பல (pala)}]{\ta{பல} (pala) — many, several}

\label{lesson:TA-C181-pala}



\begin{quote}
Before the new one: say the Tamil for favourite, then the Tamil for terrible.
\end{quote}

\subsection*{You'll want to know: \ta{பல}}
\textbf{\ta{பல}} — \emph{pala} — \textquotedblleft{}many, several\textquotedblright{}.

\begin{cousinweb}
\textbf{\ta{பல}} (\emph{pala}) is native Dravidian. DEDR 3987 lists Malayalam \emph{pala}, Kannada \emph{pala} and Telugu \emph{palu}.
\end{cousinweb}

\begin{grammarlens}[title={What you've built}]
One more word for this chapter.
\end{grammarlens}

\subsection*{Your turn}
\begin{itemize}
\raggedright
  \item \emph{Say it:} \emph{pala}
  \item \emph{Say it:} \emph{pala}, once more
  \item \emph{Say it:} say \emph{pala}
  \item \emph{Recall:} say the Tamil for favourite, then the Tamil for terrible, then say \emph{pala} again
\end{itemize}

\begin{tcolorbox}[breakable,colback=teal!4,colframe=teal!35!black,title={Before you move on}]
What does \ta{பல} mean? (\textquotedblleft{}Many, several\textquotedblright{}.) Say it once more.
\end{tcolorbox}
\section[\texorpdfstring{\ta{எல்லாம்} (ellām)}{எல்லாம் (ellām)}]{\ta{எல்லாம்} (ellām) — everything, all}

\label{lesson:TA-C181-ellam}



\begin{quote}
Before the new one: say the Tamil for terrible, then the Tamil for many.
\end{quote}

\subsection*{You'll want to know: \ta{எல்லாம்}}
\textbf{\ta{எல்லாம்}} — \emph{ellām} — \textquotedblleft{}everything, all\textquotedblright{}.

\begin{cousinweb}
\textbf{\ta{எல்லாம்}} (\emph{ellām}) is native Dravidian. DEDR 844 lists Malayalam \emph{ellām}, Kannada \emph{ella} and Telugu \emph{ella}.
\end{cousinweb}

\begin{grammarlens}[title={What you've built}]
One more word for this chapter.
\end{grammarlens}

\subsection*{Your turn}
\begin{itemize}
\raggedright
  \item \emph{Say it:} \emph{ellām}
  \item \emph{Say it:} \emph{ellām}, once more
  \item \emph{Say it:} say \emph{ellām}
  \item \emph{Recall:} say the Tamil for terrible, then the Tamil for many, then say \emph{ellām} again
\end{itemize}

\begin{tcolorbox}[breakable,colback=teal!4,colframe=teal!35!black,title={Before you move on}]
What does \ta{எல்லாம்} mean? (\textquotedblleft{}Everything, all\textquotedblright{}.) Say it once more.
\end{tcolorbox}
\section[\texorpdfstring{\ta{தடிமனான} (taṭimaṉāṉa)}{தடிமனான (taṭimaṉāṉa)}]{\ta{தடிமனான} (taṭimaṉāṉa) — thick}

\label{lesson:TA-C181-tatimanana}



\begin{quote}
Before the new one: say the Tamil for many, then the Tamil for everything.
\end{quote}

\subsection*{You'll want to know: \ta{தடிமனான}}
\textbf{\ta{தடிமனான}} — \emph{taṭimaṉāṉa} — \textquotedblleft{}thick\textquotedblright{}.

\begin{cousinweb}
\textbf{\ta{தடிமனான}} (\emph{taṭimaṉāṉa}) is \emph{taṭimaṉ}, \textquotedblleft{}thickness\textquotedblright{}, which is probably built on \emph{taṭi}, \textquotedblleft{}stick, club, a thick piece\textquotedblright{} (DEDR 3030). DEDR does not list \emph{taṭimaṉ} itself, so the link is likely rather than proven.
\end{cousinweb}

\begin{grammarlens}[title={What you've built}]
One more word for this chapter.
\end{grammarlens}

\subsection*{Your turn}
\begin{itemize}
\raggedright
  \item \emph{Say it:} \emph{taṭimaṉāṉa}
  \item \emph{Say it:} \emph{taṭimaṉāṉa}, once more
  \item \emph{Say it:} say \emph{taṭimaṉāṉa}
  \item \emph{Recall:} say the Tamil for many, then the Tamil for everything, then say \emph{taṭimaṉāṉa} again
\end{itemize}

\begin{tcolorbox}[breakable,colback=teal!4,colframe=teal!35!black,title={Before you move on}]
What does \ta{தடிமனான} mean? (\textquotedblleft{}Thick\textquotedblright{}.) Say it once more.
\end{tcolorbox}
